\chapter{The one-sided multiplicative ergodic theorem}
\label{chap:met}

This is the crux of the development. We work over a probability space \((X,\mu)\) with an
ergodic measure-preserving transformation \(T\) and a measurable cocycle generator
\(A : X \to \mathrm{Mat}_{d\times d}(\R)\) with \(\det(A x)\neq 0\), subject to the one-sided
integrability \(\log^{+}\norm{A},\,\log^{+}\norm{A^{-1}}\in L^1(\mu)\). The matrices act on
\(\R^d\) (the Euclidean space) through the operator \(v\mapsto A v\), so that the relevant norm
is the \(L^2\) operator norm, which is submultiplicative; we write \(A^{(n)}(x)\) for the cocycle
\(A(T^{n-1}x)\cdots A(x)\). The whole chapter is organized around a single positive semidefinite
limiting matrix, the \emph{Oseledets limit}
\[
  \Lambda(x) \;=\; \lim_{n\to\infty} \bigl((A^{(n)}(x))^{\top} A^{(n)}(x)\bigr)^{1/2n},
\]
whose eigenspaces, once their growth rates are matched to the deterministic singular-value
exponents, yield the filtration of the target theorem.

\section{The Oseledets limit}

The candidate approximants are the symmetric positive roots
\(q_n(x) := \bigl((A^{(n)}(x))^{\top} A^{(n)}(x)\bigr)^{1/2n}\), realized through the continuous
functional calculus on the Gram matrix. The eigenvalues of \(q_n(x)\) are the \(1/n\)-th roots of
the singular values of \(A^{(n)}(x)\); the heart of this section is that they converge a.e.\ to
the exponentials of the deterministic Lyapunov exponents.

\begin{theorem}[Deterministic singular-value exponents]\label{exists_lam_tendsto_singularValue}
  \lean{Oseledets.exists_lam_tendsto_singularValue}\leanok
  There is an antitone sequence \(\lambda^0 : \N\to\R\) (antitone on \([0,d)\)) such that, for
  every \(i<d\) and \(\mu\)-a.e.\ \(x\),
  \[
    \tfrac1n \log\bigl(\sigma_i(A^{(n)}(x))\bigr) \longrightarrow \lambda^0_i,
  \]
  where \(\sigma_i\) is the \(i\)-th sorted singular value.
\end{theorem}
\begin{proof}\leanok
  \uses{tendsto_kingman, isSubadditiveCocycle_logNorm}
  Package the ergodic limits \(\Gamma_k = \lim \tfrac1n \log s_k(A^{(n)}(x))\) of the products of
  the top \(k\) singular values, \(0\le k\le d\), obtained from Kingman's subadditive ergodic
  theorem applied to the subadditive cocycle \(\log\norm{\ext^{k}A^{(n)}}\) (the exterior-power
  functor turns submultiplicativity of the \(\ext^{k}\) operator norms into subadditivity). Then
  \(\lambda^0_i := \Gamma_{i+1}-\Gamma_i\); the difference of the two a.e.\ \(\Gamma\)-limits gives
  the \(i\)-th singular-value exponent. Antitonicity of consecutive \(\lambda^0_i\) descends from
  the antitone ordering of the singular values inside each \(n\), and chaining yields full
  antitonicity on \([0,d)\).
\end{proof}

\begin{lemma}[Block-value step function reproduces the spectrum]\label{stepVal_exp_lam}
  \lean{Oseledets.stepVal_exp_lam}\leanok
  Let \(\mathrm{stepVal}\,\lambda^0\,D\) be the step function
  \(e^{\lambda^0_{D-1}} + \sum_{k=1}^{D-1}(e^{\lambda^0_{k-1}}-e^{\lambda^0_k})\,
  \mathbf 1_{(c_k,\infty)}\) with thresholds \(c_k = e^{(\lambda^0_k+\lambda^0_{k-1})/2}\) strictly
  inside the \(k\)-th gap. If \(\lambda^0\) is antitone on \([0,D)\) and \(j<D\), then
  \(\mathrm{stepVal}\,\lambda^0\,D\,(e^{\lambda^0_j}) = e^{\lambda^0_j}\).
\end{lemma}
\begin{proof}\leanok
  At the argument \(e^{\lambda^0_j}\) the threshold indicator \(\mathbf 1_{(c_k,\infty)}\) is
  \(1\) exactly when \(k>j\) (since \(\lambda^0\) is antitone and \(c_k\) lies strictly between
  \(\lambda^0_k\) and \(\lambda^0_{k-1}\)), so only the increments above index \(j\) survive.
  Those increments \(e^{\lambda^0_{k-1}}-e^{\lambda^0_k}\) telescope to
  \(e^{\lambda^0_j}-e^{\lambda^0_{D-1}}\), which added to the constant base \(e^{\lambda^0_{D-1}}\)
  returns \(e^{\lambda^0_j}\).
\end{proof}

\begin{lemma}[Spectral deviation bound]\label{norm_sub_cfc_le_sum_eigenvalue_dev}
  \lean{Oseledets.norm_sub_cfc_le_sum_eigenvalue_dev}\leanok
  For a self-adjoint matrix \(M\) and any function \(g\),
  \[
    \norm{M - g(M)} \;\le\; \sum_j \abs{\mu_j - g(\mu_j)},
  \]
  where \(\mu_j\) ranges over the sorted eigenvalues of \(M\).
\end{lemma}
\begin{proof}\leanok
  Writing \(M = \mathrm{id}(M)\) gives \(M - g(M) = (\mathrm{id}-g)(M)\) by linearity of the
  continuous functional calculus, and the operator norm of a self-adjoint matrix functional
  calculus is the largest absolute eigenvalue \(\max_j\abs{\mu_j-g(\mu_j)}\), which is bounded by
  the full nonnegative sum.
\end{proof}

\begin{lemma}[Per-term band-projector convergence]\label{ae_forall_tendsto_block_term}
  \lean{Oseledets.ae_forall_tendsto_block_term}\leanok
  For \(\mu\)-a.e.\ \(x\) and every threshold index \(k\in[1,d)\), the block term
  \((e^{\lambda^0_{k-1}}-e^{\lambda^0_k})\cdot P^{c_k}_n(x)\) converges, where \(P^{c}_n\) is the
  band projector \(\mathbf 1_{(c,\infty)}(q_n(x))\).
\end{lemma}
\begin{proof}\leanok
  \uses{exists_lam_tendsto_singularValue}
  At a genuine gap \(\lambda^0_k < \lambda^0_{k-1}\) the threshold \(c_k\) is strictly separated
  from the two limiting eigenvalue clusters, so once the sorted eigenvalues of \(q_n(x)\) have
  converged the count of eigenvalues above \(c_k\) stabilizes and the corresponding spectral
  projector is Cauchy in operator norm. At a non-gap (\(\lambda^0_{k-1}=\lambda^0_k\)) the
  coefficient \(e^{\lambda^0_{k-1}}-e^{\lambda^0_k}\) vanishes, so the term is constantly \(0\).
\end{proof}

\begin{theorem}[Existence of the Oseledets limit]\label{tendsto_qpow}
  \lean{Oseledets.tendsto_qpow}\leanok
  For \(\mu\)-a.e.\ \(x\) the approximants \(q_n(x)\) converge in the matrix metric to a single
  matrix \(\Lambda(x)\).
\end{theorem}
\begin{proof}\leanok
  \uses{exists_lam_tendsto_singularValue,stepVal_exp_lam,norm_sub_cfc_le_sum_eigenvalue_dev,ae_forall_tendsto_block_term}
  The eigenvalues \(\mu_{j,n}=\sigma_j^{1/n}\) of \(q_n(x)\) converge a.e.\ to the exponentials
  \(e^{\lambda^0_j}\) (Theorem~\ref{exists_lam_tendsto_singularValue}). Form the block approximant
  \(\Lambda_n(x) := \mathrm{stepVal}\,\lambda^0\,d\,(q_n(x))\), a finite linear combination of band
  projectors. By Lemma~\ref{norm_sub_cfc_le_sum_eigenvalue_dev},
  \(\norm{q_n(x)-\Lambda_n(x)}\le\sum_j\abs{\mu_{j,n}-\mathrm{stepVal}(\mu_{j,n})}\); each summand
  is eventually \(\abs{\mu_{j,n}-e^{\lambda^0_j}}\to 0\) because the step function reproduces the
  exponentials on the spectrum (Lemma~\ref{stepVal_exp_lam}). Meanwhile \(\Lambda_n(x)\) converges
  as a finite sum of convergent band-projector terms
  (Lemma~\ref{ae_forall_tendsto_block_term}). Adding the two convergences gives
  \(q_n(x)\to\Lambda(x)\), and the limit is selected pointwise.
\end{proof}

\begin{definition}[The named Oseledets limit]\label{oseledetsLimit}
  \lean{Oseledets.oseledetsLimit}\leanok
  \(\Lambda(x) := (\,\lim_n (q_n(x))_{ij}\,)_{ij}\) is the entrywise real \(\limsup\)/limit of the
  matrix entries of \(q_n(x)\); it is a total, measurable function of \(x\).
\end{definition}

\begin{theorem}[The limit is the a.e.\ limit of the approximants]\label{tendsto_oseledetsLimit}
  \lean{Oseledets.tendsto_oseledetsLimit}\leanok
  For \(\mu\)-a.e.\ \(x\), \(q_n(x)\to \Lambda(x)\) in the matrix metric, and \(\Lambda\) is
  measurable.
\end{theorem}
\begin{proof}\leanok
  \uses{tendsto_qpow,oseledetsLimit}
  On the a.e.\ full convergence set of Theorem~\ref{tendsto_qpow} the entrywise limit recovers the
  matrix limit (matrix convergence in finite dimensions is entrywise), so the entrywise
  \(\liminf\) defining \(\Lambda\) equals the genuine limit. Measurability is entrywise: each
  entry is a limit of measurable functions of \(x\), and a \(\liminf\) of measurable
  \(\R\)-valued functions is measurable.
\end{proof}

\begin{proposition}[Structure of the limit]\label{oseledetsLimit_posSemidef}
  \lean{Oseledets.oseledetsLimit_posSemidef}\leanok
  For \(\mu\)-a.e.\ \(x\), \(\Lambda(x)\) is self-adjoint and positive semidefinite.
\end{proposition}
\begin{proof}\leanok
  \uses{tendsto_oseledetsLimit}
  Self-adjointness \(M^{\top}=M\) is an entrywise closed condition preserved under the matrix
  limit of the self-adjoint approximants \(q_n(x)\). For positive semidefiniteness, the quadratic
  form \(M\mapsto v^{\top}Mv\) is continuous, so
  \(v^{\top}\Lambda(x)v = \lim_n v^{\top} q_n(x) v \ge 0\) as a limit of nonnegatives.
\end{proof}

\section{The per-vector lower bound}

The lower half of the exact growth law isolates one band of the spectrum of \(q_n(x)\) and shows
that a vector with nonzero projection onto the band grows at least at the band rate.

\begin{lemma}[Gram quadratic-form band bound]\label{inner_cfc_ge_band}
  \lean{Oseledets.inner_cfc_ge_band}\leanok
  For self-adjoint \(Q\), a band indicator \(\chi=\mathbf 1_{(c,\infty)}\), and a continuous
  \(f\ge 0\) on \(\mathrm{spec}(Q)\) with \(a\le f(t)\) whenever \(c<t\),
  \[
    a\,\norm{\chi(Q)\,v}^2 \;\le\; \langle f(Q)\,v,\,v\rangle .
  \]
\end{lemma}
\begin{proof}\leanok
  The band projector \(\chi(Q)\) is a self-adjoint idempotent, so
  \(\norm{\chi(Q)v}^2 = \langle \chi(Q)v,v\rangle\). The gap operator
  \((f-a\chi)(Q)\) is positive semidefinite because \(f-a\chi\ge 0\) on the spectrum (above \(c\),
  \(f\ge a=a\chi\); below, \(\chi=0\) and \(f\ge0\)). Expanding
  \(\langle(f-a\chi)(Q)v,v\rangle\ge0\) gives the claim.
\end{proof}

\begin{lemma}[Band lower bound for the cocycle]\label{cocycle_apply_sq_ge_band}
  \lean{Oseledets.cocycle_apply_sq_ge_band}\leanok
  For \(c\ge0\) and \(n\ge1\),
  \[
    c^{2n}\,\norm{P^{c}_n(x)\,v}^2 \;\le\; \norm{A^{(n)}(x)\,v}^2 .
  \]
\end{lemma}
\begin{proof}\leanok
  \uses{inner_cfc_ge_band}
  Raising \(q_n(x)=(\mathrm{gram}_n)^{1/2n}\) to the \(2n\)-th power via the functional calculus
  recovers the Gram matrix \(\mathrm{gram}_n=(A^{(n)})^{\top}A^{(n)}\) (the composed powers compose
  to the identity on the nonnegative spectrum). Apply
  Lemma~\ref{inner_cfc_ge_band} with \(f(t)=t^{2n}\) and \(a=c^{2n}\): above \(c\) one has
  \(t^{2n}\ge c^{2n}\). The right-hand inner product is then
  \(\langle \mathrm{gram}_n v,v\rangle = \norm{A^{(n)}(x)v}^2\).
\end{proof}

\begin{lemma}[The band correction vanishes]\label{tendsto_inv_mul_log_norm_bandProjector_apply}
  \lean{Oseledets.tendsto_inv_mul_log_norm_bandProjector_apply}\leanok
  If \(P^{c}_n(x)\to P\) with \(Pv\neq0\), then
  \(\tfrac1n\log\norm{P^{c}_n(x)\,v}\to 0\).
\end{lemma}
\begin{proof}\leanok
  The evaluation \(M\mapsto Mv\) is continuous in finite dimensions, so \(P^{c}_n(x)v\to Pv\neq0\)
  and \(\norm{P^{c}_n(x)v}\to\norm{Pv}>0\). Hence the log converges to the finite number
  \(\log\norm{Pv}\), and dividing by \(n\to\infty\) sends it to \(0\).
\end{proof}

\begin{proposition}[Per-vector liminf lower bound]\label{log_le_liminf_log_cocycle_apply}
  \lean{Oseledets.log_le_liminf_log_cocycle_apply}\leanok
  If \(P^{c}_n(x)\to P\) with \(c>0\) and \(Pv\neq0\), and the cocycle growth sequence is
  cobounded, then
  \[
    \log c \;\le\; \liminf_{n}\ \tfrac1n\log\norm{A^{(n)}(x)\,v}.
  \]
\end{proposition}
\begin{proof}\leanok
  \uses{cocycle_apply_sq_ge_band,tendsto_inv_mul_log_norm_bandProjector_apply}
  Taking logs in Lemma~\ref{cocycle_apply_sq_ge_band} and dividing by \(2n\) gives, eventually,
  \[
    \log c + \tfrac1n\log\norm{P^{c}_n(x)v} \;\le\; \tfrac1n\log\norm{A^{(n)}(x)v}.
  \]
  The left side converges to \(\log c\) since the band-correction term vanishes
  (Lemma~\ref{tendsto_inv_mul_log_norm_bandProjector_apply}). Passing to the \(\liminf\) along the
  inequality, using its boundedness and the coboundedness of the right side (supplied by the
  Furstenberg--Kesten integrability of the top exponent), yields the bound.
\end{proof}

\begin{lemma}[Band-projector nesting, kernel propagation]\label{limitBandProjector_apply_eq_zero_of_le}
  \lean{Oseledets.limitBandProjector_apply_eq_zero_of_le}\leanok
  For thresholds \(c\le c'\) with limit band projectors \(P,P'\), if \(Pv=0\) then \(P'v=0\).
\end{lemma}
\begin{proof}\leanok
  The finite-\(n\) bands are nested: \(\mathbf 1_{(c,\infty)}\cdot\mathbf 1_{(c',\infty)}=
  \mathbf 1_{(c',\infty)}\) on the spectrum since \((c',\infty)\subseteq(c,\infty)\), so
  \(P^{c}_n P^{c'}_n = P^{c'}_n\); passing to the limit gives \(PP'=P'\). Both limit projectors are
  symmetric (limits of self-adjoint matrices), so transposing gives \(P'P=P'\), whence
  \(P'v = P'(Pv) = 0\).
\end{proof}

\section{The spectral upper bound and the determinant squeeze}

The upper half is the genuinely non-elementary step. For a vector \(v\) in the slow subspace of
\(\Lambda\) at level \(e^t\), one bounds \(\norm{A^{(n)}(x)v}\) by the restricted operator norm of
\(A^{(n)}\) on the slow subspace, whose growth exponent is pinned --- non-circularly --- by a
determinant/volume squeeze.

\begin{lemma}[Tempering]\label{tempering_posLog}
  \lean{Oseledets.tempering_posLog}\leanok
  If \(T\) is measure-preserving and \(g\ge0\) with \(g\in L^1(\mu)\), then for \(\mu\)-a.e.\ \(x\),
  \(\tfrac1n\,g(T^n x)\to 0\).
\end{lemma}
\begin{proof}\leanok
  The series \(\sum_n g(T^n x)/n^{2}\) is a.e.\ finite by integrability and invariance of \(\mu\),
  so its terms \(g(T^n x)/n^2\to0\); the Borel--Cantelli/sublinear-growth argument then gives
  \(g(T^n x)=o(n)\), i.e.\ \(\tfrac1n g(T^n x)\to0\).
\end{proof}

\begin{lemma}[Slow-volume exponent squeeze]\label{limsup_topSlow_le_of_squeeze}
  \lean{Oseledets.limsup_topSlow_le_of_squeeze}\leanok
  Given a volume cocycle whose top, slow and remaining log-exponent sequences satisfy a sum law
  \(\mathrm{vol} = \mathrm{slow} + \mathrm{rest}\) with the appropriate limits, the slow
  restricted-operator-norm exponent obeys
  \(\limsup_n \tfrac1n\log(\mathrm{slow}_n) \le \lambda_i\).
\end{lemma}
\begin{proof}\leanok
  The total volume exponent is the Furstenberg--Kesten determinant limit \(\sum_j\lambda^0_j\); the
  fast-block volume is the exterior-power Kingman limit; their difference forces the slow-block
  volume exponent. The squeeze converts a sum identity of limits into an upper bound for the slow
  factor once the fast and remaining factors are pinned, the angle/tilt between fast and slow
  blocks tempering to zero so no cross term inflates the slow volume.
\end{proof}

\begin{theorem}[Spectral upper bound via the determinant squeeze]\label{spectral_upper_bound_of_squeeze}
  \lean{Oseledets.spectral_upper_bound_of_squeeze}\leanok
  Let \(v\neq0\) lie in the slow subspace. Given the slow restricted-norm exponent bound
  \(\limsup_n \tfrac1n\log r_n\le \lambda_i\) and the restriction estimate
  \(\norm{A^{(n)}(x)v}\le r_n\norm v\) (valid for \(v\) slow), one has
  \[
    \limsup_{n}\ \tfrac1n\log\norm{A^{(n)}(x)\,v} \;\le\; \lambda_i.
  \]
\end{theorem}
\begin{proof}\leanok
  \uses{limsup_topSlow_le_of_squeeze,tempering_posLog}
  Take logs in \(\norm{A^{(n)}(x)v}\le r_n\norm v\), divide by \(n\), and pass to the \(\limsup\):
  the \(\norm v\) factor contributes \(0\) and \(r_n\) contributes
  \(\limsup_n\tfrac1n\log r_n\le\lambda_i\) (Lemma~\ref{limsup_topSlow_le_of_squeeze}). The
  restricted-norm exponent \(r_n\) depends only on the \emph{global} volume cocycle and the
  tempered slow--fast angle, never on the growth of \(v\) itself, so the argument is
  non-circular.
\end{proof}

\begin{theorem}[Upper bound on the slow space]\label{vslow_subset_lambdaSublevel_of_upper}
  \lean{Oseledets.vslow_subset_lambdaSublevel_of_upper}\leanok
  On the ultrametric-growth good set, every vector \(v\) of the \(\Lambda\)-slow band
  \(\mathrm{vslow}(e^t)\) with \(\limsup_n\tfrac1n\log\norm{A^{(n)}v}\le t\) lies in the growth
  sublevel \(\{v:\overline\lambda(x,v)\le t\}\).
\end{theorem}
\begin{proof}\leanok
  \uses{spectral_upper_bound_of_squeeze,lambdaSublevel}
  The \(\limsup\) of \(\tfrac1n\log\norm{A^{(n)}v}\) is by definition the upper growth function
  \(\overline\lambda(x,v)\); the hypothesis says it is \(\le t\), which is exactly membership in
  the sublevel \(\mathrm{lambdaSublevel}\,t\). The zero vector lies in every submodule.
\end{proof}

\section{Spectral identification of the filtration}

The two bounds match the spectral filtration of \(\Lambda\) with the analytic limsup filtration.
The bridge is that the finite-\(n\) band projectors converge a.e.\ to the functional calculus
indicator of \(\Lambda\).

\begin{theorem}[Band projectors converge to the CFC indicator]\label{ae_tendsto_bandProjector_cfc_indicator}
  \lean{Oseledets.ae_tendsto_bandProjector_cfc_indicator}\leanok
  For \(\mu\)-a.e.\ \(x\), every \(c>0\) that is not one of the limiting eigenvalues
  \(e^{\lambda_{\mathrm{sing}}(x,i)}\) of \(\Lambda(x)\) satisfies
  \[
    P^{c}_n(x) \longrightarrow \mathbf 1_{(c,\infty)}(\Lambda(x)).
  \]
\end{theorem}
\begin{proof}\leanok
  \uses{tendsto_oseledetsLimit,oseledetsLimit_posSemidef}
  Since \(c\) avoids the spectrum of \(\Lambda(x)\), there is a gap \(\delta>0\) between \(c\) and
  every eigenvalue. Replace the discontinuous indicator by a continuous \(\delta/2\)-clamp
  surrogate \(\chi\) that is Lipschitz and agrees with \(\mathbf 1_{(c,\infty)}\) at distance
  \(\ge\delta/2\) from \(c\). Once the sorted eigenvalues of \(q_n(x)\) are within \(\delta/2\) of
  their limits, \(\chi\) and \(\mathbf 1_{(c,\infty)}\) agree on both spectra, so
  \(P^{c}_n(x)=\chi(q_n(x))\) eventually; by Lipschitz continuity of the functional calculus and
  \(q_n(x)\to\Lambda(x)\) (Theorem~\ref{tendsto_oseledetsLimit}),
  \(\chi(q_n(x))\to\chi(\Lambda(x))=\mathbf 1_{(c,\infty)}(\Lambda(x))\).
\end{proof}

\begin{theorem}[Reverse slow-flag inclusion]\label{ae_lambdaSublevel_le_vslow}
  \lean{Oseledets.ae_lambdaSublevel_le_vslow}\leanok
  For \(\mu\)-a.e.\ \(x\) and every \(t\),
  \(\mathrm{lambdaSublevel}(x,t) \le \mathrm{vslow}(x,e^t)\).
\end{theorem}
\begin{proof}\leanok
  \uses{ae_tendsto_bandProjector_cfc_indicator,log_le_liminf_log_cocycle_apply,limitBandProjector_apply_eq_zero_of_le}
  Contrapositively, a vector \(v\notin\mathrm{vslow}(e^t)\) has nonzero component in the band of
  \(\Lambda(x)\) above \(e^t\); by Theorem~\ref{ae_tendsto_bandProjector_cfc_indicator} the finite
  band projectors converge to the corresponding CFC indicator with \(Pv\neq0\), so the per-vector
  lower bound (Proposition~\ref{log_le_liminf_log_cocycle_apply}) forces
  \(\liminf\tfrac1n\log\norm{A^{(n)}v} > t\), hence \(\overline\lambda(x,v)>t\) and
  \(v\notin\mathrm{lambdaSublevel}\,t\). Kernel propagation across nested thresholds
  (Lemma~\ref{limitBandProjector_apply_eq_zero_of_le}) makes this consistent for all \(t\).
\end{proof}

\begin{theorem}[The slow flag equals the limsup sublevel]\label{vslow_eq_lambdaSublevel_of_upper}
  \lean{Oseledets.vslow_eq_lambdaSublevel_of_upper}\leanok
  Under the spectral upper bound and the reverse inclusion, for \(\mu\)-a.e.\ \(x\) and every
  \(t\),
  \[
    \mathrm{vslow}(x,e^t) \;=\; \mathrm{lambdaSublevel}(x,t).
  \]
\end{theorem}
\begin{proof}\leanok
  \uses{vslow_subset_lambdaSublevel_of_upper,ae_lambdaSublevel_le_vslow}
  Two inclusions: the forward one
  (Theorem~\ref{vslow_subset_lambdaSublevel_of_upper}, from the upper bound) shows slow vectors
  grow slowly, hence lie in the sublevel; the reverse one
  (Theorem~\ref{ae_lambdaSublevel_le_vslow}, from the lower bound) shows slowly-growing vectors lie
  in the slow space. Antisymmetry gives the identity, simultaneously for all \(t\) on the a.e.\
  ultrametric-growth good set.
\end{proof}

\section{Ruelle's reverse cofactor bound and the top-gap envelope}

The upper bound on an individual vector reduces, after diagonalizing the limit, to controlling the
overlap matrix between the sorted Gram eigenbasis at level \(n\) and the limiting eigenbasis. The
graded overlap is controlled by a leakage induction; Ruelle's cofactor estimate then converts a
one-sided forward decay into the full pairwise rate.

\begin{lemma}[Ruelle reverse cofactor bound]\label{entry_reverse_bound_of_orthogonal}
  \lean{Oseledets.RuelleCofactor.entry_reverse_bound_of_orthogonal}\leanok
  Let \(S\) be orthogonal (\(S S^{\top}=1\)) with the graded forward decay
  \(\abs{S_{ab}}\le c\cdot e^{-\max(g_b-g_a,0)}\). Then every entry obeys the reverse bound at the
  full pairwise rate:
  \[
    \abs{S_{ij}} \;\le\; (d-1)!\,c^{\,d-1}\,e^{-(g_i-g_j)} .
  \]
\end{lemma}
\begin{proof}\leanok
  Since \(S^{-1}=S^{\top}\), the entry \(S_{ij}\) is \((\det S)^{-1}\) times the cofactor
  \(\mathrm{adj}(S)_{ji}\), and \(\abs{\det S}=1\) because \(SS^{\top}=1\). Expanding the minor by
  the Leibniz formula, every surviving permutation term collects the level imbalance \(g_i-g_j\)
  by telescoping the forward factors, each forward factor contributing at most
  \(c\,e^{-\max(\cdot,0)}\); there are at most \((d-1)!\) such terms.
\end{proof}

\begin{definition}[Top-gap fast-band-mass envelope]\label{TopGapMassEnvelope}
  \lean{Oseledets.TopGapMassEnvelope}\leanok
  \(\mathrm{TopGapMassEnvelope}\,A\,T\,\lambda^0\,x\) asserts the uniform geometric leakage of
  fast-band mass across each genuine gap: for every cut tolerance \(\delta\) there is a constant
  \(C\) controlling, eventually and uniformly, the band mass that crosses the top gap of each
  stratum.
\end{definition}

\begin{lemma}[Multi-source geometric envelope]\label{multi_source_envelope}
  \lean{Oseledets.ChainCore.multi_source_envelope}\leanok
  If a nonnegative sequence \(a\) obeys a one-step recursion fed by finitely many source sequences
  each decaying geometrically with ratio \(\rho<1\), then \(a_n\) is bounded by a fixed multiple of
  the summed source envelopes for all \(n\).
\end{lemma}
\begin{proof}\leanok
  Each single source contributes a geometric partial sum bounded by \(K/(1-\rho)\); summing the
  finitely many per-source envelopes and folding them through the linear one-step recursion gives a
  uniform bound on the chained quantity \(a_n\).
\end{proof}

\begin{lemma}[Per-stratum envelope step]\label{perStratumEnvelope_step}
  \lean{Oseledets.perStratumEnvelope_step}\leanok
  At a fixed cut strictly inside a gap of width \(\ge G\), the one-step band-mass increment is the
  current mass damped by \(e^{-G}\) plus a tempered source term; iterating produces the per-stratum
  leakage envelope.
\end{lemma}
\begin{proof}\leanok
  \uses{multi_source_envelope}
  In the sorted-Gram-eigenbasis block decomposition, the band mass above the cut at step \(n+1\)
  is the mass at step \(n\) attenuated by the singular-value ratio across the gap (bounded by
  \(e^{-nG}\)-type damping after \(n\) iterates) plus the contribution injected by the one-step
  generator \(A(T^n x)\), whose log-norm is tempered to \(o(n)\). This is exactly the geometric
  one-step recursion fed by tempered sources of Lemma~\ref{multi_source_envelope}.
\end{proof}

\begin{theorem}[The top-gap envelope, a.e.]\label{topGapMassEnvelope_ae}
  \lean{Oseledets.topGapMassEnvelope_ae}\leanok
  For \(\mu\)-a.e.\ \(x\), the top-gap fast-band-mass envelope
  \(\mathrm{TopGapMassEnvelope}\,A\,T\,\lambda^0\,x\) holds.
\end{theorem}
\begin{proof}\leanok
  \uses{TopGapMassEnvelope,perStratumEnvelope_step,exists_lam_tendsto_singularValue,tempering_posLog}
  Fix the deterministic distinct gap \(G>0\) separating distinct exponents. On the a.e.\ set where
  every singular-value exponent converges (Theorem~\ref{exists_lam_tendsto_singularValue}) and the
  one-step generator log-norm is tempered (Lemma~\ref{tempering_posLog}), build the per-stratum
  leakage envelope for each gap pair (Lemma~\ref{perStratumEnvelope_step}) and assemble them into
  the top-gap envelope: each stratum's fast-band mass crossing its top gap is geometrically
  damped, uniformly in the cut, by a single per-pair constant maximized over the finitely many
  pairs.
\end{proof}

\section{Constancy of the spectrum}

The deterministic exponent set is constant in \(x\) by construction, so once the per-point spectrum
is identified with it, ergodic constancy is automatic.

\begin{lemma}[Spectrum identity from two inclusions]\label{lyapunovSpectrum_eq_of_subsets}
  \lean{Oseledets.lyapunovSpectrum_eq_of_subsets}\leanok
  If at \(x\) every realized exponent is a deterministic one and every deterministic exponent is
  attained, then \(\mathrm{lyapunovSpectrum}(x) = \mathrm{distinctExp}\,\lambda^0\,d\).
\end{lemma}
\begin{proof}\leanok
  \uses{lyapunovSpectrum}
  Both directions are finite-set inclusions; antisymmetry of \(\subseteq\) gives the equality of
  the two finite subsets of \(\R\).
\end{proof}

\begin{theorem}[Ergodic constancy of the spectrum]\label{lyapunovSpectrum_eq_distinctExp_of_lambdaBar}
  \lean{Oseledets.lyapunovSpectrum_eq_distinctExp_of_lambdaBar}\leanok
  Given a.e.\ that every realized value of the upper growth function is a deterministic exponent
  (upper inclusion) and every deterministic exponent is attained (lower inclusion), for
  \(\mu\)-a.e.\ \(x\),
  \[
    \mathrm{lyapunovSpectrum}(x) \;=\; \mathrm{distinctExp}\,\lambda^0\,d .
  \]
\end{theorem}
\begin{proof}\leanok
  \uses{lyapunovSpectrum_eq_of_subsets,vslow_eq_lambdaSublevel_of_upper,lyapunovSpectrum}
  On the ultrametric-growth good set the two Finset inclusions are equivalent to native statements
  about the upper growth function \(\overline\lambda(x,\cdot)\): the upper inclusion is the
  spectral-upper-bound output (each stratum value is a deterministic exponent), the lower inclusion
  is attainment from the lower bound. Lemma~\ref{lyapunovSpectrum_eq_of_subsets} then gives the
  identity with the deterministic constant set; since that set does not depend on \(x\), the
  identification is \(T\)-invariant and the spectrum is a.e.\ constant --- ergodicity adds nothing
  beyond the deterministic value already in hand.
\end{proof}

\section{Assembling the target theorem}

The forward graded-overlap bound (built from the top-gap envelope and Ruelle's reverse bound) gives
the spectral upper bound \(\mathrm{hupper}\); combined with the reverse inclusion it yields the slow
flag, and the lower bound supplies attainment. Together they discharge the three a.e.\ interfaces
(spectrum, slow flag, exact growth) of the filtration assembly.

\begin{theorem}[Filtration from the spectral upper bound]\label{oseledets_filtration_of_upper}
  \lean{Oseledets.oseledets_filtration_of_upper}\leanok
  Assume the per-vector spectral upper bound on the slow flag, the reverse slow-flag inclusion, the
  two spectrum inclusions, and the band-projector convergence datum. Then there exist
  \(k\), strictly decreasing \(\lambda:\mathrm{Fin}\,k\to\R\), and a measurable family
  \(V\) forming a.e.\ a strictly decreasing \(A\)-equivariant flag along which
  \(\tfrac1n\log\norm{A^{(n)}(x)v}\to\lambda_i\) on each stratum.
\end{theorem}
\begin{proof}\leanok
  \uses{vslow_eq_lambdaSublevel_of_upper,lyapunovSpectrum_eq_distinctExp_of_lambdaBar,log_le_liminf_log_cocycle_apply,vflag,lambdaBar}
  The deterministic exponents \(\lambda^0\) come from
  Theorem~\ref{exists_lam_tendsto_singularValue}. The spectrum interface is discharged by the two
  inclusions (constancy, Theorem~\ref{lyapunovSpectrum_eq_distinctExp_of_lambdaBar}); the slow-flag
  interface by \(\mathrm{vslow}=\mathrm{lambdaSublevel}\)
  (Theorem~\ref{vslow_eq_lambdaSublevel_of_upper}); the exact-growth interface by combining the
  unconditional upper half (the stratum value \(\overline\lambda=\lambda_i\) on
  \(\mathrm{vflag}\)), the lower half (Proposition~\ref{log_le_liminf_log_cocycle_apply} fed the
  band datum), and Furstenberg--Kesten boundedness into a two-sided limit. These three interfaces
  feed the generic slow-flag assembly.
\end{proof}

\begin{theorem}[Filtration from the top-gap envelope]\label{oseledets_filtration_of_topgap}
  \lean{Oseledets.oseledets_filtration_of_topgap}\leanok
  Under the standing ergodic, invertible, log-integrable hypotheses, and assuming the top-gap
  envelope \(\mathrm{TopGapMassEnvelope}\) quantified over \(\lambda^0\), the full Oseledets
  filtration conclusion holds.
\end{theorem}
\begin{proof}\leanok
  \uses{oseledets_filtration_of_upper,topGapMassEnvelope_ae,entry_reverse_bound_of_orthogonal,vslow_eq_lambdaSublevel_of_upper,ae_tendsto_bandProjector_cfc_indicator,TopGapMassEnvelope}
  Diagonalize \(\Lambda(x)\) by its limit eigenbasis with eigenvalues \(e^{\lambda_{\mathrm{sing}}}\)
  and slow-orthogonality. The forward graded-overlap bound, consuming the envelope
  (Definition~\ref{TopGapMassEnvelope}), yields the one-sided forward decay of the overlap matrix
  between the sorted Gram eigenbasis and the limit eigenbasis; Ruelle's reverse cofactor estimate
  (Lemma~\ref{entry_reverse_bound_of_orthogonal}) upgrades this to the full pairwise rate, which is
  exactly the slow-restriction bound feeding the spectral upper bound \(\mathrm{hupper}\) on the
  limit slow space. The band-projector convergence
  (Theorem~\ref{ae_tendsto_bandProjector_cfc_indicator}) supplies the reverse slow-flag inclusion
  and the lower-bound datum, and the spectrum inclusions follow from the slow-flag identity. These
  discharge all hypotheses of Theorem~\ref{oseledets_filtration_of_upper}.
\end{proof}

\begin{theorem}[One-sided Oseledets multiplicative ergodic theorem]\label{oseledets_filtration}
  \lean{Oseledets.oseledets_filtration}\leanok
  Let \(\mu\) be a probability measure, \(T:X\to X\) ergodic measure-preserving, and
  \(A:X\to\mathrm{Mat}_{d\times d}(\R)\) measurable with \(\det(A x)\neq 0\) and
  \(\log^{+}\norm{A},\,\log^{+}\norm{A^{-1}}\in L^1(\mu)\). Then there are \(k\) distinct Lyapunov
  exponents \(\lambda:\mathrm{Fin}\,k\to\R\), strictly decreasing, and a measurable family of
  subspaces
  \[
    V : \mathrm{Fin}(k+1)\to X\to \mathrm{Submodule}_{\R}\bigl(\R^d\bigr)
  \]
  with each \(x\mapsto V_i\,x\) measurable, such that for \(\mu\)-a.e.\ \(x\):
  \(V_0\,x=\top\), \(V_k\,x=\bot\); the flag is strictly decreasing,
  \(V_{i+1}\,x < V_i\,x\); it is \(A\)-equivariant,
  \(A(x)\,V_i\,x = V_i\,(Tx)\); and along it the cocycle grows at the exact rate \(\lambda_i\):
  \[
    \tfrac1n\log\norm{A^{(n)}(x)\,v}\ \longrightarrow\ \lambda_i
    \qquad\text{for every } v\in V_i\,x\setminus V_{i+1}\,x .
  \]
\end{theorem}
\begin{proof}\leanok
  \uses{oseledets_filtration_of_topgap,topGapMassEnvelope_ae}
  If \(d=0\) the trivial flag \(\top=\bot\) with no exponents discharges the statement. For
  \(d>0\), the top-gap envelope holds a.e.\ (Theorem~\ref{topGapMassEnvelope_ae}), so the
  conditional assembly Theorem~\ref{oseledets_filtration_of_topgap} applies directly and produces
  the exponents, the measurable equivariant flag, and the exact per-stratum growth limits.
\end{proof}
